\section{Secret Sharing --- Shamir}
%%% generated by the blueprint pipeline from TCSlib.Cryptography.SecretSharing.Shamir
\subsection*{Secret Sharing --- Shamir: Introduction}

Declarations of the Lean module \texttt{TCSlib.Cryptography.SecretSharing.Shamir}.

\subsection*{Secret Sharing --- Shamir: Results}

\begin{definition}[Nonzero field elements as parties]
\lean{TCSlib.Shamir.Party}
\statementsource{Sha79}{pdf-a76120fbec44-p002-b003}
\statementsource{BGW88}{pdf-b3f55f2a4332-p003-b005}
\statementsource{AL17}{pdf-27f2b3787521-p004-b002}
\label{TCSlib.Shamir.Party}
\leanok
For a field $F$, a party is a nonzero element of $F$.
\end{definition}

\begin{definition}[Party label]
\lean{TCSlib.Shamir.label}
\statementsource{Sha79}{pdf-a76120fbec44-p002-b003}
\statementsource{BGW88}{pdf-b3f55f2a4332-p003-b005}
\statementsource{AL17}{pdf-27f2b3787521-p004-b002}
\label{TCSlib.Shamir.label}
\leanok
\uses{TCSlib.Shamir.Party}
For a field $F$, the label of a party is the nonzero field element that defines that
party.
\end{definition}

\begin{theorem}[Injectivity of party labels]
\lean{TCSlib.Shamir.label_injective}
\label{TCSlib.Shamir.label_injective}
\leanok
\uses{TCSlib.Shamir.label}
\difficulty{1}
Let $F$ be a field. The map assigning to each nonzero element of $F$ its label is
injective.
\end{theorem}

\begin{definition}[Secret polynomial]
\lean{TCSlib.Shamir.poly}
\statementsource{Sha79}{pdf-a76120fbec44-p002-b003}
\statementsource{BGW88}{pdf-b3f55f2a4332-p003-b005}
\statementsource{AL17}{pdf-27f2b3787521-p004-b002}
\label{TCSlib.Shamir.poly}
\leanok
Let $F$ be a field, let $k$ be a natural number, let $s\in F$, and let $q$ be a
polynomial of degree less than $\max(k-1,0)$. The polynomial associated to $s$ and $q$
is $s+Xq$.
\end{definition}

\begin{theorem}[Secret at Zero]
\lean{TCSlib.Shamir.poly_eval_zero}
\proofsource{Sha79}{pdf-a76120fbec44-p002-b003}
\statementsource{BGW88}{pdf-b3f55f2a4332-p003-b005}
\label{TCSlib.Shamir.poly_eval_zero}
\leanok
\proofstep
  {TCSlib.SecretSharing.AccessStructure}
  {context}
  {Sha79}
  {pdf-a76120fbec44-p001-b005}
\proofstep
  {TCSlib.SecretSharing.AccessStructure.threshold}
  {direct}
  {Sha79}
  {pdf-a76120fbec44-p001-b005}
\proofstep
  {TCSlib.SecretSharing.AccessStructure.threshold_auth}
  {direct}
  {Sha79}
  {pdf-a76120fbec44-p001-b005}
\proofstep
  {TCSlib.Shamir.poly}
  {context}
  {Sha79}
  {pdf-a76120fbec44-p002-b003}
\uses{TCSlib.Shamir.poly}
\difficulty{1}
Let $F$ be a field, let $k$ be a natural number, let $s\in F$, and let $q$ be a
polynomial of degree less than $k-1$. The polynomial $s+Xq(X)$ evaluates to $s$ at zero.
\end{theorem}

\begin{theorem}[Evaluation of the Sharing Polynomial]
\lean{TCSlib.Shamir.poly_eval}
\statementsource{Sha79}{pdf-a76120fbec44-p002-b003}
\label{TCSlib.Shamir.poly_eval}
\leanok
\uses{TCSlib.Shamir.poly}
\difficulty{1}
Let $F$ be a field, let $k$ be a natural number, let $s\in F$, and let $q$ be a
polynomial of degree less than $k-1$. For every $a\in F$, the polynomial $s+Xq(X)$
evaluates to $s+a q(a)$.
\end{theorem}

\begin{theorem}[Degree bound for the sharing polynomial]
\lean{TCSlib.Shamir.degree_poly_lt}
\label{TCSlib.Shamir.degree_poly_lt}
\leanok
\uses{TCSlib.Shamir.poly}
\difficulty{5}
Let $F$ be a field, let $k>0$, let $\mathrm{sec}\in F$, and let $q$ be a polynomial of
degree less than $k-1$. The polynomial $\mathrm{sec}+Xq$ has degree less than $k$.
\end{theorem}

\begin{definition}[Polynomial secret-sharing scheme]
\lean{TCSlib.Shamir.scheme}
\statementsource{Sha79}{
  pdf-a76120fbec44-p002-b002,
  pdf-a76120fbec44-p002-b003
}
\statementsource{BGW88}{
  pdf-b3f55f2a4332-p003-b005,
  pdf-b3f55f2a4332-p003-b006
}
\statementsource{AL17}{pdf-27f2b3787521-p004-b002}
\label{TCSlib.Shamir.scheme}
\leanok
\uses{TCSlib.SecretSharing.Scheme, TCSlib.Shamir.Party, TCSlib.Shamir.label, TCSlib.Shamir.poly}
Let $F$ be a field and $k$ a natural number. The scheme has nonzero elements of $F$ as
parties, and both secrets and shares are elements of $F$; its randomness consists of
polynomials of degree less than $\max(k-1,0)$. Given a secret $s$, randomness $q$, and
party $a$, the share is the value at $a$ of the polynomial $s+Xq$. For a finite set of
parties and a share assignment on all parties, reconstruction interpolates the shares on
that set and evaluates the interpolating polynomial at $0$.
\end{definition}

\begin{theorem}[Share evaluation]
\lean{TCSlib.Shamir.scheme_share}
\proofsource{BGW88}{pdf-b3f55f2a4332-p003-b005}
\label{TCSlib.Shamir.scheme_share}
\leanok
\proofstep
  {TCSlib.SecretSharing.AccessStructure}
  {context}
  {BGW88}
  {pdf-b3f55f2a4332-p003-b005}
\proofstep
  {TCSlib.SecretSharing.AccessStructure.threshold}
  {context}
  {BGW88}
  {pdf-b3f55f2a4332-p003-b005}
\proofstep
  {TCSlib.SecretSharing.AccessStructure.threshold_auth}
  {context}
  {BGW88}
  {pdf-b3f55f2a4332-p003-b005}
\proofstep
  {TCSlib.SecretSharing.Scheme}
  {context}
  {BGW88}
  {pdf-b3f55f2a4332-p003-b005,
    pdf-b3f55f2a4332-p003-b006}
\proofstep
  {TCSlib.Shamir.Party}
  {context}
  {BGW88}
  {pdf-b3f55f2a4332-p003-b005}
\proofstep
  {TCSlib.Shamir.label}
  {context}
  {BGW88}
  {pdf-b3f55f2a4332-p003-b005}
\proofstep
  {TCSlib.Shamir.poly}
  {direct}
  {BGW88}
  {pdf-b3f55f2a4332-p003-b005}
\proofstep
  {TCSlib.Shamir.poly_eval_zero}
  {context}
  {BGW88}
  {pdf-b3f55f2a4332-p003-b005,
    pdf-b3f55f2a4332-p003-b006}
\proofstep
  {TCSlib.Shamir.scheme}
  {context}
  {BGW88}
  {pdf-b3f55f2a4332-p003-b005,
    pdf-b3f55f2a4332-p003-b006}
\uses{TCSlib.Shamir.Party, TCSlib.Shamir.label, TCSlib.Shamir.poly, TCSlib.Shamir.scheme}
\difficulty{1}
Let $F$ be a field, let $k$ be a natural number, let $s\in F$, let $q$ be a polynomial
of degree less than $k-1$, and let $a\in F$ be nonzero. In the specified sharing scheme,
the share of $s$ using randomness $q$ for party $a$ is the value at $a$ of the
polynomial $s+Xq(X)$.
\end{theorem}

\begin{theorem}[Locality of Shamir reconstruction]
\lean{TCSlib.Shamir.localReconstruct}
\label{TCSlib.Shamir.localReconstruct}
\leanok
\uses{TCSlib.SecretSharing.Scheme.LocalReconstruct, TCSlib.Shamir.scheme}
\difficulty{3}
Let $F$ be a field and let $k$ be a nonnegative integer. In the Shamir scheme over $F$
with parameter $k$, reconstruction from a finite set $A$ of parties depends only on the
shares held by parties in $A$: whenever two share functions agree on every party in $A$,
reconstruction from $A$ gives the same result for both.
\end{theorem}

\begin{theorem}[Correctness of Shamir sharing]
\lean{TCSlib.Shamir.correct}
\proofsource{AL17}{pdf-27f2b3787521-p004-b002}
\label{TCSlib.Shamir.correct}
\leanok
\proofstep
  {TCSlib.SecretSharing.AccessStructure}
  {context}
  {AL17}
  {pdf-27f2b3787521-p028-b003}
\proofstep
  {TCSlib.SecretSharing.AccessStructure.threshold}
  {context}
  {AL17}
  {pdf-27f2b3787521-p028-b003}
\proofstep
  {TCSlib.SecretSharing.AccessStructure.threshold_auth}
  {context}
  {AL17}
  {pdf-27f2b3787521-p028-b003}
\proofstep
  {TCSlib.SecretSharing.Scheme}
  {context}
  {AL17}
  {pdf-27f2b3787521-p004-b002}
\proofstep
  {TCSlib.SecretSharing.Scheme.Correct}
  {direct}
  {AL17}
  {pdf-27f2b3787521-p004-b002}
\proofstep
  {TCSlib.Shamir.Party}
  {context}
  {AL17}
  {pdf-27f2b3787521-p004-b002}
\proofstep
  {TCSlib.Shamir.label}
  {context}
  {AL17}
  {pdf-27f2b3787521-p004-b002}
\proofstep
  {TCSlib.Shamir.label_injective}
  {context}
  {AL17}
  {pdf-27f2b3787521-p004-b002}
\proofstep
  {TCSlib.Shamir.poly}
  {direct}
  {AL17}
  {pdf-27f2b3787521-p004-b002}
\proofstep
  {TCSlib.Shamir.poly_eval_zero}
  {direct}
  {AL17}
  {pdf-27f2b3787521-p004-b002}
\proofstep
  {TCSlib.Shamir.degree_poly_lt}
  {direct}
  {AL17}
  {pdf-27f2b3787521-p004-b002}
\proofstep
  {TCSlib.Shamir.scheme}
  {direct}
  {AL17}
  {pdf-27f2b3787521-p004-b002}
\uses{TCSlib.SecretSharing.AccessStructure.threshold, TCSlib.SecretSharing.Scheme.Correct, TCSlib.Shamir.Party, TCSlib.Shamir.degree_poly_lt, TCSlib.Shamir.label, TCSlib.Shamir.label_injective, TCSlib.Shamir.poly, TCSlib.Shamir.poly_eval_zero, TCSlib.Shamir.scheme}
\difficulty{4}
Let $F$ be a field and let $k>0$. In the Shamir scheme over $F$, parties are the nonzero
elements of $F$, and sharing a secret $\mathrm{sec}$ with randomness $q$ gives each
party the value at its label of the polynomial $\mathrm{sec}+Xq$, where $q$ has degree
less than $k-1$. Reconstruction from a set $A$ of parties evaluates at zero the
polynomial interpolating their shares; for every set $A$ with at least $k$ elements,
every secret, and every randomness value, this reconstruction returns the secret.
\end{theorem}

\begin{theorem}[Perfect Privacy of Threshold Sharing]
\lean{TCSlib.Shamir.perfectPrivacy}
\statementsource{Sha79}{pdf-a76120fbec44-p002-b004}
\statementsource{BGW88}{
  pdf-b3f55f2a4332-p003-b005,
  pdf-b3f55f2a4332-p003-b006
}
\statementsource{AL17}{
  pdf-27f2b3787521-p013-b002,
  pdf-27f2b3787521-p013-b003,
  pdf-27f2b3787521-p014-b002
}
\label{TCSlib.Shamir.perfectPrivacy}
\leanok
\uses{TCSlib.SecretSharing.AccessStructure.threshold, TCSlib.SecretSharing.Scheme.PerfectPrivacy, TCSlib.Shamir.Party, TCSlib.Shamir.label, TCSlib.Shamir.label_injective, TCSlib.Shamir.poly_eval, TCSlib.Shamir.scheme, TCSlib.Shamir.scheme_share}
\difficulty{5}
Let $F$ be a field, let $k$ be a natural number, and consider the scheme that shares a
secret $s\in F$ by evaluating $s+Xq(X)$ at nonzero field elements, where $q$ has degree
less than $k-1$; reconstruction from a set of parties evaluates the interpolating
polynomial at zero. For every finite set $B$ of nonzero field elements with $|B|<k$ and
every pair of secrets $s,s'\in F$, there is a bijection $e$ of the randomness space such
that, for every randomness choice $r$ and every $b\in B$, the share of $s'$ using
randomness $e(r)$ at $b$ equals the share of $s$ using randomness $r$ at $b$.
\end{theorem}

\begin{theorem}[Number of Compatible Sharing Polynomials]
\lean{TCSlib.Shamir.natCard_coins_view}
\statementsource{BGW88}{
  pdf-b3f55f2a4332-p003-b005,
  pdf-b3f55f2a4332-p003-b006
}
\proofsource{AL17}{
  pdf-27f2b3787521-p013-b002,
  pdf-27f2b3787521-p013-b003
}
\label{TCSlib.Shamir.natCard_coins_view}
\leanok
\proofstep
  {TCSlib.SecretSharing.AccessStructure}
  {context}
  {AL17}
  {pdf-27f2b3787521-p028-b003}
\proofstep
  {TCSlib.SecretSharing.AccessStructure.threshold}
  {context}
  {AL17}
  {pdf-27f2b3787521-p004-b002}
\proofstep
  {TCSlib.SecretSharing.AccessStructure.threshold_auth}
  {context}
  {AL17}
  {pdf-27f2b3787521-p004-b002}
\proofstep
  {TCSlib.SecretSharing.Scheme}
  {context}
  {AL17}
  {pdf-27f2b3787521-p004-b002}
\proofstep
  {TCSlib.Interpolation.fixedEvalEquiv}
  {context}
  {AL17}
  {pdf-27f2b3787521-p004-b002,
    pdf-27f2b3787521-p015-b004}
\proofstep
  {TCSlib.Interpolation.natCard_degreeLT_eval_eq}
  {context}
  {AL17}
  {pdf-27f2b3787521-p004-b002}
\proofstep
  {TCSlib.Shamir.Party}
  {context}
  {AL17}
  {pdf-27f2b3787521-p004-b002}
\proofstep
  {TCSlib.Shamir.label}
  {context}
  {AL17}
  {pdf-27f2b3787521-p004-b002}
\proofstep
  {TCSlib.Shamir.label_injective}
  {context}
  {AL17}
  {pdf-27f2b3787521-p004-b002}
\proofstep
  {TCSlib.Shamir.poly}
  {direct}
  {AL17}
  {pdf-27f2b3787521-p004-b002}
\proofstep
  {TCSlib.Shamir.poly_eval_zero}
  {direct}
  {AL17}
  {pdf-27f2b3787521-p004-b002}
\proofstep
  {TCSlib.Shamir.poly_eval}
  {context}
  {AL17}
  {pdf-27f2b3787521-p004-b002}
\proofstep
  {TCSlib.Shamir.scheme}
  {direct}
  {AL17}
  {pdf-27f2b3787521-p004-b002}
\proofstep
  {TCSlib.Shamir.scheme_share}
  {direct}
  {AL17}
  {pdf-27f2b3787521-p004-b002}
\uses{TCSlib.Interpolation.natCard_degreeLT_eval_eq, TCSlib.Shamir.Party, TCSlib.Shamir.label, TCSlib.Shamir.label_injective, TCSlib.Shamir.poly_eval, TCSlib.Shamir.scheme, TCSlib.Shamir.scheme_share}
\difficulty{5}
Let $F$ be a finite field and let $k>0$ satisfy $k\le |F|$. For a finite set $B$ of
nonzero elements of $F$ with $|B|<k$, a secret $s\in F$, and any prescribed values
$w(b)\in F$ for $b\in B$, the number of polynomials of the form $s+Xq(X)$ with $\deg
q<k-1$ whose evaluations at every $b\in B$ equal $w(b)$ is $|F|^{k-1-|B|}$.
\end{theorem}
